\section{Introduction}


\subsection{The Capital Efficiency Problem}

Traditional DeFi lending protocols like MakerDAO and Compound require overcollateralization ratios of 150-200\%, resulting in significant capital inefficiency. A user depositing \$100k can only borrow \$50-66k, locking substantial value unproductively.

Self-repaying loan protocols like Alchemix improved capital efficiency by using yield to automatically repay debt, but maintained conservative 50\% LTV ratios and limited collateral types to yield-bearing assets on Ethereum only.

\subsection{Lux Credit Innovation}

Lux Credit, launched in December 2023, addresses these limitations through:

\begin{enumerate}
\item \textbf{90\% LTV Ratios}: Highest capital efficiency in DeFi through automated yield repayment
\item \textbf{11\% APY on LUX}: Sustainable yield from diversified strategies
\item \textbf{Cross-Chain Collateral}: Bitcoin, Ethereum, and 15+ assets via M-Chain bridge
\item \textbf{MPC Security}: Threshold custody eliminates centralized key management
\item \textbf{Zero Liquidations}: Conservative risk model since December 2023 launch
\item \textbf{Self-Repaying}: Automated debt repayment from yield generation
\end{enumerate}

\subsection{Key Achievements (December 2023 - Q3 2024)}

\begin{table}[h]
\centering
\begin{tabular}{@{}lr@{}}
\toprule
\textbf{Metric} & \textbf{Value} \\ \midrule
Loans Processed & 18,400 \\
Total Value Locked & \$427M \\
Average LTV & 87.3\% \\
Liquidations & 0 \\
Yield Generated & \$31.2M \\
Active Users & 12,847 \\
Average Loan Duration & 387 days \\
Protocol Revenue & \$2.1M \\ \bottomrule
\end{tabular}
\caption{Lux Credit performance metrics (Q3 2024)}
\end{table}

